\documentclass[11pt]{article}
\usepackage[T1]{fontenc}
\usepackage{lmodern}
\usepackage[margin=1in]{geometry}
\usepackage{amsmath,amssymb,amsthm}
\newcommand{\mathscr}{\mathcal}
\numberwithin{equation}{section}
\usepackage[expansion=false]{microtype}
\usepackage[hidelinks,hypertexnames=false]{hyperref}
\setlength{\emergencystretch}{2em}
\newcommand{\workbenchurl}{https://github.com/KokunoYumeto/zeta-and-l-functions-workbench}
\newcommand{\recordurl}{https://github.com/KokunoYumeto/mathematics-workbenches}
\newcommand{\F}{\mathcal F}
\newcommand{\B}{\mathcal B}
\newcommand{\HH}{\mathcal H}
\newcommand{\N}{\mathrm N}
\newcommand{\cont}{\operatorname{cont}}
\newcommand{\Rea}{\operatorname{Re}}
\newcommand{\Ima}{\operatorname{Im}}
\newtheorem{theorem}{Theorem}[section]
\newtheorem{corollary}[theorem]{Corollary}
\newtheorem{lemma}[theorem]{Lemma}
\title{Poles and residue subtraction for a cubic Hecke content series,\\after Hecke}
\author{}
\date{24 November 2026}
\begin{document}
\maketitle
\begin{abstract}
The zero-count estimates of Bennett, Martin, O'Bryant, and Rechnitzer
(2021) and Hasanalizade, Shen, and Wong (2022) supply the explicit
pole-count constants in this paper. Liu's stated zero-free theorem (2026) is a separate external input for the numerical example. For a fixed cubic Hecke character over the
Eisenstein field, we extract every Euler term linear in the primitive
norm variable by an alternating product of shifted Hecke functions.
The remaining Euler product converges normally whenever the norm
variable has real part greater than one half and the content variable
has positive real part. We prove that the rational-content quotient
has infinitely many poles on or to the right of the quarter line,
even after every possible cancellation by numerator zeros. All these
poles survive in the two-variable series at the single norm parameter
$s=4$. Under zeta RH they lie on the quarter line. This obstructs a
full holomorphic tube; it does not establish a stronger bound for
the cubic prime sum or prove its Riemann hypothesis.
We also construct an absolutely convergent expansion in complete content-residue blocks. It includes all multiplicities, the ordinary cofactor, and every finite norm-cap collision. At the stated zero-free boundary the bulk and complete single residues are below square-root size for content cutoff exponent \(0.28\), but the available complete double term has exponent \(0.5065405965\ldots\). We identify its exact missing estimate and prove no smaller support cutoff. Standard Hecke analytic facts, the two published zero-count estimates, and Liu's zero-free statement are explicitly granted external inputs. The mathematics and text
were produced by GPT-6 Astra in ChatGPT Pro mode, with a separate statement-only
reconstruction and targeted checks; no human has checked this work.
\end{abstract}

\section{Definitions and full statements}

Let $K=\mathbb Q(\omega)$, where $\omega^2+\omega+1=0$, and let
$\psi$ be either faithful cubic Hecke character for
$K(\sqrt[3]{269})/K$. We take as fixed input that its conductor is
$(269)$ and that $\psi((d))=1$ for rational integers $d$ prime to
$269$. Delete all prime ideals over $3$ and $269$; this deletion is
denoted by $S$. Every ideal sum below is restricted to ideals prime
to $S$. For a nonzero integral ideal $\mathfrak a$, let
$\cont(\mathfrak a)$ be the greatest positive integer $d$ such that
$\mathfrak a\subseteq d\mathcal O_K$. Put
\begin{equation}\label{eq:defs}
 P(s)=L_K^S(s,\psi),\qquad
 m(\mathfrak a)=\mu_K(\mathfrak a)\psi(\mathfrak a),\qquad
 r(\mathfrak a)=\cont(\mathfrak a).
\end{equation}
For $\Rea s>1$ and $\Rea v>1$, define
\begin{equation}\label{eq:F}
 \F(s,v)=\sum_{\mathfrak a}
 m(\mathfrak a)(\N\mathfrak a)^{-s}
 r(\mathfrak a)^{2s-v}.
\end{equation}
Let $\eta=\chi_{-3}$ and set
\begin{equation}\label{eq:A}
 A(v)=\frac{L(v,\eta)}
 {\zeta(2v)(1-3^{-2v})(1-269^{-v})}.
\end{equation}
Complex powers of positive numbers use their real logarithms.

\paragraph{Explicit analytic grants.}
We use continuation, functional equations, and fixed-strip polynomial
bounds for these fixed Hecke functions; see Zaman (2016, \S2.1,
equations (2.1)--(2.3) and Lemmas 2.1--2.3, preprint pp.~6--7).
The reciprocal bound on closed inner zero-free strips follows from
Lemma~\ref{am:reciprocal-growth} below and the absolutely convergent
Euler product. The published Hecke continuation and strip-growth
inputs are not independently re-proved here.
The field and character data just stated are also fixed inputs.
The two explicit published zero-count estimates are stated exactly
in Section~\ref{sec:count}; their analytic proofs are not reconstructed.
The zero-free hypothesis in Theorem~\ref{thm:mixed} is an assumption,
not a new zero-free theorem.

For an undeleted rational prime put $x=p^{-s}$ and $y=p^{-v}$.
If $p\equiv1\pmod3$, choose one prime ideal $\mathfrak p$ over $p$
and put $\alpha_p=\psi(\mathfrak p)$ and
$a_p=\alpha_p+\alpha_p^{-1}\in\{2,-1\}$. Define
\begin{equation}\label{eq:Dh}
\begin{array}{c|c|c}
 &D_p(z)&h_p(s,v)\\ \hline
p\equiv1\pmod3&1-a_pz+z^2&1-a_px/(1+y)\\[2pt]
p\equiv2\pmod3&1-z^2&1
\end{array}
\end{equation}
There are no local factors at $p=3$ or $269$.

\begin{theorem}[All mixed factors]\label{thm:mixed}
Suppose $1/2<\beta\le1$ and $P$ is zero-free on $\Rea s>\beta$.
Then
\begin{equation}\label{eq:B}
 \B(s,v)=\prod_{j\ge1}P(s+jv)^{(-1)^{j+1}}
\end{equation}
is a nowhere-zero holomorphic function on
$\Rea s>\beta$, $\Rea v>0$. Its tail converges normally in logarithms,
so its value is unchanged by reordering its factors. Define
\begin{equation}\label{eq:Elocal}
 E_{\infty,p}(s,v)=
 \frac{h_p(s,v)}{D_p(x)}
 \prod_{j\ge1}D_p(xy^j)^{(-1)^{j+1}}.
\end{equation}
The local $j$-products converge absolutely in logarithms, and
\begin{equation}\label{eq:E}
 E_\infty(s,v)=\prod_{p\ne3,269}E_{\infty,p}(s,v)
\end{equation}
converges normally on $\Rea s>1/2$, $\Rea v>0$, independently of
prime order when the local factors are grouped as in
\eqref{eq:Elocal}. It may have zeros. There is a meromorphic
continuation of \eqref{eq:F} to $\Rea s>\beta$, $\Rea v>0$ given by
\begin{equation}\label{eq:factor}
 \F(s,v)=A(v)\HH(s,v),\qquad
 \HH(s,v)=P(s)^{-1}\B(s,v)E_\infty(s,v).
\end{equation}
The function $\HH$ is holomorphic there, may have zeros, and grows
at most polynomially in $1+|\Ima s|+|\Ima v|$ on every bounded
closed inner real-part strip. Moreover,
\begin{equation}\label{eq:limit}
 \HH(s,v)\longrightarrow1\quad(s\longrightarrow+\infty,
 \ s\in\mathbb R)
\end{equation}
uniformly on every half-plane $\Rea v\ge\nu_0>0$.
\end{theorem}

\begin{theorem}[Survival of poles]\label{thm:poles}
If $A$ has a pole of order $k$ at $v_0$ with $\Rea v_0>0$, then,
for every sufficiently large real $s$, $\F(s,v)$ has a pole of
exactly order $k$ at $v_0$. More strongly,
\begin{equation}\label{eq:s4}
 \HH(4,v)\ne0\qquad(\Rea v\ge1/4).
\end{equation}
Every pole of $A$ in this closed half-plane therefore survives in
$\F(4,v)$ with its full order. The fixed-$s$ assertion uses only
the ordinary Euler half-plane and requires no additional zero-free
hypothesis on $P$.
\end{theorem}

Let $N_\zeta^+(T)$ count nontrivial zeta zeros with
$0<\Ima\rho\le T$, with multiplicity. Let $N_\eta^{\pm}(T)$ count
zeros of $L(s,\eta)$ in $0<\Rea s<1$, $|\Ima s|\le T$, again with
multiplicity. Let $Q_A(T)$ be the sum of pole orders of $A$ in
\begin{equation}\label{eq:region}
 1/4\le\Rea v<1/2,\qquad 0<\Ima v\le T.
\end{equation}

\begin{theorem}[Cancellation-inclusive count]\label{thm:count}
For every $T>0$,
\begin{equation}\label{eq:count}
 Q_A(T)\ge\frac12N_\zeta^+(2T)-\frac14N_\eta^{\pm}(T).
\end{equation}
This allows all possible common zeros and all multiplicities.
Define
\begin{align}
 E_\zeta(t)&=0.1038\log t+0.2573\log\log t+9.3675,\label{eq:Ez}\\
 E_\eta(t)&=\min\{0.247\log(3t)+6.894,\,
                      0.298\log(3t)+4.358\}.\label{eq:Eeta}
\end{align}
For $T\ge e/2$ the published estimates give
\begin{equation}\label{eq:explicit}
 Q_A(T)\ge\frac{T}{4\pi}\log\frac{2T}{3\pi e}
          -\frac12E_\zeta(2T)-\frac14E_\eta(T).
\end{equation}
Consequently
\begin{equation}\label{eq:asymp}
 Q_A(T)\ge\frac{T}{4\pi}\log\frac{2T}{3\pi e}-O(\log T).
\end{equation}
\end{theorem}

\begin{corollary}[Exact method boundary]\label{cor:barrier}
For $u_0<1/4$, the series \eqref{eq:F} has no continuation
holomorphic on the entire tube $\Rea s>\beta$, $\Rea v>u_0$.
If the Riemann hypothesis for zeta is assumed, $A$ is holomorphic
on $\Rea v>1/4$, but has infinitely many poles on its boundary,
all of which survive in $\F(4,v)$. This is only a barrier to a full
holomorphic tube. It does not exclude subtracting residues,
cancellation between residue contributions, a different transform,
or a stronger estimate for the actual balanced prime sum. It gives
no stronger bound for the cubic balanced prime sum or new Riemann-hypothesis theorem.
\end{corollary}


\subsection{Content tails and complete residue terms}

For the residue assertions assume, in addition, that
\(1/2<\beta<1\), \(1/2<\rho<1\), and that \(\zeta\) has no zero
in \(\Rea s>\rho\); its simple pole at \(1\) is allowed. Set
\[
 u=\rho/2,\qquad \ell=(1-\rho)/2,\qquad
 G(s,v)=\B(s,v)E_\infty(s,v).
\]
Let \(b,\chi,\chi_d:[0,\infty)\to[0,1]\) be fixed smooth
functions equal to \(1\) on \([0,1]\) and \(0\) on \([2,\infty)\),
and put \(K=1-\chi\). For \(V\in C_c^\infty((1,2))\),
\(Y\ge2\), \(R\ge1\), \(U^2\ge Y\), and \(q=U^2/Y\), define
\begin{equation}\label{res:tail-def}
 J_V(Y,U,R)=\sum_{\mathfrak a_1,\mathfrak a_2}
 m(\mathfrak a_1)m(\mathfrak a_2)
 b(\N\mathfrak a_1/U)b(\N\mathfrak a_2/U)
 V(\N\mathfrak a_1\N\mathfrak a_2/Y)
 K(r(\mathfrak a_1)r(\mathfrak a_2)/R).
\end{equation}
There is no cross-coprimality condition. Put
\[
 \Psi_q(t_1,t_2)=b(t_1)b(t_2)V(qt_1t_2),\qquad
 \widehat\Psi_q(s_1,s_2)=
 \int_{(0,\infty)^2}\Psi_q(t_1,t_2)t_1^{s_1}t_2^{s_2}
 \frac{dt_1dt_2}{t_1t_2},
\]
and \(\widetilde K(z)=\int_0^\infty K(t)t^{-z}\,dt/t\)
for \(\Rea z>0\).
All uniformity assertions below require a common compact support
for \(V\) and bounds on finitely many logarithmic derivatives.

\begin{theorem}[Complete content-residue subtraction]\label{res:main}
Fix
\[
 0<v_L<\ell<u<v_R<1/2,\qquad
 \sigma=\beta+\delta,\qquad \alpha=2\sigma-v_L<3,
 \qquad \delta>0.
\]
Every positive-real-part pole of \(A\), after numerator cancellation,
has real part in \([\ell,u]\). There exist
\(T_N\in[N+1/4,N+3/4]\) such that \(A\) is polynomially bounded
on both vertical lines and the horizontal segments at \(\pm T_N\).
Partition the intervening strip into the central rectangle and
the positive and negative rectangles between successive heights.
For each such rectangle \(Q_N\) define
\[
 \mathcal R_N[\Phi]=\frac1{2\pi i}
 \int_{\partial Q_N}^{+}A(v)\Phi(v)\,dv.
\]
For holomorphic \(\Phi\) near the rectangle this is its full
residue sum, with all multiplicities and cancellations, and
\[
 |\mathcal R_N[\Phi]|
 \ll(1+|N|)^C\sup_{\partial Q_N}|\Phi|.
\]
Set
\begin{align*}
 \mathcal K_z(v_1,v_2)
 ={}&\widetilde K(z)R^{-z}U^{z+(v_1+v_2)/2}\\
 &{}\cdot\widehat\Psi_q((z+v_1)/2,(z+v_2)/2)
 \HH((z+v_1)/2,v_1)\HH((z+v_2)/2,v_2).
\end{align*}
Then \(J_V=I_L+S_1+S_2+D\), where
\begin{align*}
 I_L={}&\frac1{4(2\pi i)^3}
 \int_{(\alpha)}\int_{(v_L)}\int_{(v_L)}
 A(v_1)A(v_2)\mathcal K_z(v_1,v_2)\,dv_2\,dv_1\,dz,\\
 S_1={}&\frac1{4(2\pi i)^2}
 \int_{(\alpha)}\int_{(v_L)}
 A(v_2)\sum_N\mathcal R_N^{(1)}
 [\mathcal K_z(\cdot,v_2)]\,dv_2\,dz,\\
 D={}&\frac1{4(2\pi i)}
 \int_{(\alpha)}\sum_{N,M}
 \mathcal R_N^{(1)}\mathcal R_M^{(2)}[\mathcal K_z]\,dz,
\end{align*}
and \(S_2\) is symmetric. These integrals and the sums of complete
blocks are absolutely convergent. No absolute convergence of
individual zero residues is asserted. With \(\Delta=v_R-v_L\),
\begin{align*}
 |I_L|&\ll Y^\sigma R^{-\alpha}\log(2+q),\\
 |S_1|+|S_2|&\ll
 Y^{\sigma+\Delta/4}q^{\Delta/4}R^{-\alpha}\log(2+q),\\
 |D|&\ll
 Y^{\sigma+\Delta/2}q^{\Delta/2}R^{-\alpha}\log(2+q).
\end{align*}
For arbitrary \(q\ge1\) one also has
\begin{equation}\label{res:main-general-q}
 |D|\ll_\varepsilon\log(2+q)
 \left\{Y^{\beta+\varepsilon}R^{-(2\beta-u)}
 +Y^\sigma R^{-\alpha}\bigl[1+(Yq)^{\Delta/4}\bigr]\right\}.
\end{equation}
In the balanced regime \(q=Y^{o(1)}\), this implies
\begin{equation}\label{res:main-balanced}
 D\ll_\varepsilon Y^{\beta+\varepsilon}
 R^{-(2\beta-u)}\log(2+q)
 \qquad (R\ge Y^{1/4}).
\end{equation}
The hypothesis on \(q\) in this last assertion is essential to
the deduction from the stated estimates.
\end{theorem}

\begin{theorem}[Cofactors and finite norm caps]\label{res:caps-main}
Let
\[
 a(\mathfrak c)=-\tfrac12\psi(\mathfrak c)
 d_{2,K}(\mathfrak c)\log\N\mathfrak c,\qquad
 M_U(s,z)=\sum_{\mathfrak d}m(\mathfrak d)b(\N\mathfrak d/U)
 (\N\mathfrak d)^{-s}r(\mathfrak d)^z.
\]
The function \(M_U\) is entire. Without an \(n\)-dependent direct
weight, the full ordinary cofactor for a global content cutoff is
\(P(s)P'(s)\), and its uncapped simplification is
\[
 P(s)P'(s)\F(s,2s-z)^2
 =A(2s-z)^2G(s,2s-z)^2\,\frac{P'(s)}{P(s)}.
\]
Zeros of \(G\) may cancel the possible logarithmic-derivative poles.
The cutoff \(R_{\mathfrak c}=(X/\N\mathfrak c)^\gamma\) changes
the cofactor argument to \(s-\gamma z\).
An actual direct cutoff retains its weighted finite cofactor sum.

The cap transform has the continuation
\[
 \widehat b(w)=\frac1w+\int_1^2b(t)t^w\frac{dt}{t},
\]
and, initially in an absolute domain,
\[
 M_U(s,z)=\frac1{2\pi i}\int_{(c)}
 \widehat b(w)U^w\F(s+w,2(s+w)-z)\,dw.
\]
At a content pole \(v_0\), the moving cap pole is
\(w=(z+v_0)/2-s\). Its collision with \(w=0\) cancels the entire
artificial principal part for every pole order. The four corresponding
double-cap combinations cancel their complete principal parts as
well. At \(U^2=2X\), the two moving residues retain the exact power
\[
 (X/n)^{s-\gamma z}U^{z+(v_1+v_2)/2-2s}
 =
 2^{-s+z/2+(v_1+v_2)/4}
 n^{-(s-\gamma z)}
 X^{(1/2-\gamma)z+(v_1+v_2)/4}.
\]
Thus a pole of an uncapped term is not a genuine pole of the finite
capped sum, and canceling that artificial pole supplies no omitted
power saving.
\end{theorem}

\subsection{Balanced specialization and the remaining estimate}

Let \(L_X=X^{o(1)}\), \(U=\sqrt{2X}\), \(n=\N\mathfrak c\),
\(Y=X/n\), and
\[
 V_n(t)=W(t)\bigl[1-\chi_d((L_X/n)t)\bigr],\qquad
 \|(t\,d/dt)^jW\|_\infty
 \le CH^j j!\{\log(e+j)\}^{rj},\quad r>1.
\]
Here \(W\) is fixed and supported compactly in \((1,2)\).
The weight is zero for \(n\ge2L_X\); its required finite seminorms
are uniform in \(n\). Define the complete weighted functional
\[
 \mathscr D_\gamma(X)=
 \sum_{\N\mathfrak c<2L_X}a(\mathfrak c)
 D_{V_n}(X/n,\sqrt{2X},(X/n)^\gamma),
\]
and define \(\mathscr I_\gamma,\mathscr S_{i,\gamma},
\mathscr J_\gamma\) by the same weighting.

\paragraph{External zero-free input for the numerical example.}
We take Liu's stated theorem for finite-order unitary Hecke
characters as an explicit external input
\cite[Theorem~1.2, p.~3; Theorem~8.29, pp.~78--80]{AMLiu}.
Its specialization to the present character gives \(\beta=B_*\),
and its trivial-character specialization over \(\mathbb Q\)
gives \(\rho=B_*\), with the pole of zeta allowed.
The zero-free theorem's analytic proof is not reconstructed here.
The numerical calculation uses the rational enclosure
\[
 0.873345856106718576801184796065<B_*<
 0.873345856106718576801184796066,
\]
for the constant specified in the cited equations~(42), (44).
It adds no zeta RH hypothesis.

\begin{corollary}[Exact remaining double-residue target]\label{res:numerical-main}
Under that explicit input, take \(v_L=1/20\), \(\gamma=7/25\).
For every \(\varepsilon>0\),
\begin{align*}
 \mathscr I_\gamma&\ll_\varepsilon
 X^{(11/25)B_*+7/500+\varepsilon},\\
 |\mathscr S_{1,\gamma}|+|\mathscr S_{2,\gamma}|
 &\ll_\varepsilon X^{(113/200)B_*+3/2000+\varepsilon},\\
 \mathscr D_\gamma&\ll_\varepsilon X^{(29/50)B_*+\varepsilon}.
\end{align*}
The first two exponents are strictly below \(1/2\), whereas
\((29/50)B_*=0.5065405965418967745\ldots>1/2\).
Consequently the stronger cutoff \(0.28\), compared with
the previous exponent
\[ (2B_*-1)/(3B_*)=0.2849927501963953104\ldots, \]
would require
exactly the additional estimate
\[
 \mathscr D_{7/25}(X)\ll_\varepsilon X^{1/2+\varepsilon}.
\]
It is equivalent to that square-root target for
\(\mathscr J_{7/25}\), since the other three terms already have
a fixed margin. The required saving over the available bound is
\[ (29/50)B_*-1/2=0.0065405965418967745\ldots. \]
No such saving or smaller support cutoff is proved here.
\end{corollary}


\section{Proof overview}

At a split rational prime, the squarefree ideal sum has local
factor $1-a_px+y$. Its content-only part is $1+y$; at an inert
prime the whole factor is already content-only. Dividing by these
content factors leaves $1-a_px/(1+y)$. The alternating shifted
Hecke factors reproduce exactly the geometric series
$y/(1+y)$ that is needed to remove the remaining linear term in
$x$. After that cancellation every term in the remainder has
at least $x^2$, which gives normal prime convergence for
$\Rea s>1/2$.

On an inner tube, only finitely many shifted Hecke factors lie
outside an ordinary Euler half-plane. Those finitely many factors
are controlled by the zero-free hypothesis and the reciprocal-growth
lemma, and the remaining factors form a uniformly convergent
geometric tail. The direct split-prime product then proves both
the limit as $s$ tends to infinity and nonvanishing at $s=4$.

For the pole count, reflection puts at least half of the
positive-height zeta zeros on or to the right of its critical line.
Every numerator zero capable of cancellation lies strictly left
of the Dirichlet critical line and belongs to a four-element
symmetry orbit. The exact counting conventions of the two cited
papers therefore lead to different coefficients, $1/2$ and $1/4$.
This argument includes all multiplicities and does not assume
that zeta and $L(s,\eta)$ have disjoint zero sets.

For residue subtraction, an averaged logarithmic minimum-modulus argument gives polynomially bounded good heights without zero spacing. Smooth Mellin decay makes the complete block sums absolutely convergent. The ordinary cofactor reduces reciprocal multiplicity, while the finite-cap expansion cancels every artificial local principal part. Its compensating powers retain the complete double-residue cost, isolating the additional estimate needed for a smaller cutoff.

\section{The squarefree local product and its finite corrections}

At a split prime $p=\mathfrak p\overline{\mathfrak p}$, the
squarefree choices are $1,\mathfrak p,\overline{\mathfrak p},(p)$.
Their norms are $1,p,p,p^2$ and their contents are $1,1,1,p$.
Since $\psi((p))=1$, the character values at the two primes are
$\alpha_p$ and $\alpha_p^{-1}$. Thus the local sum is
\begin{equation}\label{eq:splitF}
 1-(\alpha_p+\alpha_p^{-1})p^{-s}
       +p^{-2s}p^{2s-v}=1-a_px+y.
\end{equation}
At an inert prime, $(p)$ itself is prime of norm $p^2$ and content
$p$, with character value $1$. Its squarefree sum is $1-y$.
For $\sigma=\Rea s>1$ and $\nu=\Rea v>1$, the sum of the absolute
values of the nonconstant local terms is bounded by
\begin{equation}
 \sum_{p\text{ split}}(2p^{-\sigma}+p^{-\nu})
       +\sum_{p\text{ inert}}p^{-\nu}<\infty.
\end{equation}
Multiplicativity across distinct rational prime supports therefore
justifies the absolutely convergent Euler product for \eqref{eq:F}.

At an undeleted prime the local factor of $L(v,\eta)/\zeta(2v)$ is
\begin{equation}
 \frac{1-y^2}{1-\eta(p)y}=1+\eta(p)y.
\end{equation}
This is $1+y$ at a split prime and $1-y$ at an inert prime.
At $3$, ramification of $\eta$ leaves the unwanted factor
$1-3^{-2v}$; the first correction in \eqref{eq:A} removes it.
The prime $269\equiv2\pmod3$ is inert, and the second correction
removes its factor $1-269^{-v}$. Thus
\begin{equation}\label{eq:Aproduct}
 A(v)=\prod_{p\text{ split}}(1+p^{-v})
       \prod_{p\text{ inert},\ p\ne269}(1-p^{-v})
       \qquad(\Rea v>1).
\end{equation}
No factor at $3$ or $269$ remains.

It follows initially that $\F=A\HH_0$, where
\begin{equation}\label{eq:H0}
 \HH_0(s,v)=\prod_{p\text{ split}}
       \left(1-\frac{a_pp^{-s}}{1+p^{-v}}\right)
       =\prod_{p\text{ split}}h_p(s,v).
\end{equation}
This product is normally convergent on $\sigma>1$, $\nu>0$:
for $\sigma\ge1+\epsilon$, $\nu\ge\nu_0>0$,
\begin{equation}\label{eq:Hbound}
 \sum_{p\text{ split}}|h_p-1|
 \le\frac{2}{1-2^{-\nu_0}}\sum_p p^{-1-\epsilon}<\infty.
\end{equation}
The denominators $1+p^{-v}$ are nonzero there.
Consequently $A\HH_0$ already supplies meromorphic continuation
to $\sigma>1$, $\nu>0$, with no extra zero-free hypothesis.

\begin{lemma}[Reciprocal growth in a zero-free half-plane]
\label{am:reciprocal-growth}
Let \(f\) be holomorphic and nonvanishing in \(\Re z>a\), with
polynomial growth on every bounded closed vertical strip there.
Suppose that on one fixed line \(\Re z=b>a\),
\(|f(b+it)|\ge c>0\) for all real \(t\). Then \(1/f\) has
polynomial growth on every bounded closed strip strictly inside
\(\Re z>a\).
\end{lemma}
\begin{proof}
Fix a target real-part interval \([a+\delta,B]\) and put
\(T=3+|t|\). Choose a fixed-width rectangle around the horizontal
segments joining \(b+it\) to this interval at height \(t\),
with fixed positive margins and left edge to the right of \(a\).
Polynomial strip growth gives
\(\log|f(z)|\le C_0+A\log T\) throughout that rectangle.
Increase \(C_0\) so that
\[
 h_t(z)=C_0+A\log T-\log|f(z)|
\]
is positive and harmonic there. Harmonicity follows from
nonvanishing; no global choice of a complex logarithm is needed.
Cover each horizontal segment by a uniformly bounded number of
disks of fixed positive radius, with successive centers at most
half a radius apart. The Poisson formula for a positive harmonic
function bounds its value within half a disk radius by three times
its value at the center. Successive applications therefore give
\[
 h_t(\sigma+it)\le C_Hh_t(b+it)\ll\log T,
 \qquad a+\delta\le\sigma\le B.
\]
The last inequality uses the assumed lower bound at \(b+it\).
Thus \(\log|f(\sigma+it)|\ge-C\log T\), proving
\(|1/f(\sigma+it)|\ll T^C\), with constants independent of
\(t\) and \(\sigma\) in the target strip.
\end{proof}

For the present function $P$, take $a=\beta$ and $b=2$.
The stated Hecke theory gives polynomial strip growth. On
$\Rea s=2$, the absolutely convergent reciprocal Euler product
is uniformly bounded, so $|P(2+it)|$ is bounded below by a
positive constant. The lemma therefore proves polynomial growth
of $1/P$ on every closed inner zero-free strip, without an
assumption about zero multiplicities or a separate zero-count input.

\section{Complete mixed extraction}

\subsection{The shifted product of continued Hecke functions}

Assume the hypothesis of Theorem~\ref{thm:mixed}. On a compact
subset of $\sigma>\beta$, $\nu>0$, choose $\delta,\nu_0>0$ and
an integer $J$ such that
$\Rea(s+jv)\ge1+\delta$ for $j>J$. In the Euler half-plane the
canonical absolutely convergent logarithm satisfies
\begin{equation}\label{eq:logP}
 |\log P(w)|\le C_\delta 2^{-\Rea w}
       \qquad(\Rea w\ge1+\delta).
\end{equation}
Indeed, the logarithmic prime-ideal series has bounded character
coefficients. Bounding the number of ideals of a given norm by
the quadratic divisor bound reduces it to a convergent Dirichlet
series; factoring out $2^{-\Rea w}$ gives a constant depending
only on $\delta$ and the fixed field. The bound is uniform in
$\Ima w$.

The series of logarithms for $j>J$ is therefore dominated by a
geometric series in $2^{-j\nu_0}$, locally uniformly in $(s,v)$.
Its exponential is holomorphic and nowhere zero. The first $J$
factors in \eqref{eq:B} are also holomorphic and nowhere zero:
$P$ is holomorphic by the finite-order Hecke grant and nonzero
by the explicit hypothesis. Their reciprocals introduce no poles.
This proves the existence and nonvanishing of $\B$. Absolute
convergence of the tail logarithms also proves independence of
the order of the $j$-factors.

\subsection{The exact local remainder and prime convergence}

Every root of $D_p$ in \eqref{eq:Dh} has modulus $1$.
Hence $D_p(xy^j)$ has no zero for $\sigma>0$, $\nu>0$ and $j\ge0$.
For each fixed prime, its logarithm is given by the power series
at $1$. The $j$-product in \eqref{eq:Elocal} converges normally
and absolutely in logarithms: its tail is
$O_p(|x||y|^j)$ at split primes and $O_p(|x|^2|y|^{2j})$ at inert
primes. Thus every $E_{\infty,p}$ is holomorphic on
$\sigma>1/2$, $\nu>0$. Its numerator $h_p$ may vanish; it is
never inverted.

Fix $\sigma\ge\sigma_0>1/2$ and $\nu\ge\nu_0>0$.
For all sufficiently large primes, use the uniformly convergent
power-series logarithms
\begin{align}
 \log D_p(z)&=-a_pz+O(|z|^2) &&(p\text{ split}),\label{eq:Dexp}\\
 \log D_p(z)&=O(|z|^2) &&(p\text{ inert}).
\end{align}
At a split prime,
\begin{equation}
 \log h_p=-\frac{a_px}{1+y}
          +O\left(\frac{|x|^2}{|1+y|^2}\right).
\end{equation}
The linear term of $\log E_{\infty,p}$ is exactly
\begin{align}
 -\frac{a_px}{1+y}+a_px
 -a_px\sum_{j\ge1}(-1)^{j+1}y^j
 &=-\frac{a_px}{1+y}+a_px-\frac{a_pxy}{1+y}\notag\\
 &=0.\label{eq:cancel}
\end{align}
The remaining terms are $O_{\nu_0}(|x|^2)$, because
$|1+y|\ge1-2^{-\nu_0}$ and
$\sum_{j\ge1}|y|^{2j}$ is uniformly bounded. At an inert prime
all terms are already quadratic. Exponentiating yields
\begin{equation}\label{eq:Eestimate}
 E_{\infty,p}(s,v)=1+O_{\sigma_0,\nu_0}(p^{-2\sigma})
\end{equation}
for all sufficiently large $p$, uniformly in both imaginary parts.
Since $\sum_p p^{-2\sigma_0}<\infty$, this proves normal convergence
of \eqref{eq:E}. The finitely many smaller local factors are
holomorphic and uniformly bounded on any fixed inner real-part
strip: their denominators have positive lower bounds because
$|x|<1$ and $|y|<1$. Thus $E_\infty$ is uniformly bounded there
in both imaginary variables.

The product is independent of prime order with each local
cancellation kept grouped. This does not license separating
linearly divergent prime subproducts outside their ordinary
Euler domains. Zeros of a local $h_p$ are permitted, so the
conclusion is holomorphy of $E_\infty$, not nonvanishing.

\subsection{Identification with the original series and growth}

For $\sigma>1$, $\nu>1$, all prime and $j$ logarithmic sums used
here are absolutely convergent. The local factor of $P(s)$ is
$D_p(x)^{-1}$, and the local factor of $\B$ is
\begin{equation}
 \prod_{j\ge1}D_p(xy^j)^{(-1)^j}.
\end{equation}
Multiplying this with $D_p(x)$ and \eqref{eq:Elocal} leaves exactly
$h_p$. Fubini's theorem therefore gives
$\HH_0=P(s)^{-1}\B E_\infty$ in the initial absolute domain.
The right side is holomorphic on $\sigma>\beta$, $\nu>0$ and
hence supplies the continuation in \eqref{eq:factor}. In the larger
domain $\B$ is its $j$-product of continued Hecke functions;
we have not assumed an unjustified prime Euler product there.

For growth, fix an inner strip
\begin{equation}\label{eq:inner}
 \beta+\epsilon\le\sigma\le M,\qquad
 \nu_0\le\nu\le N,
\end{equation}
with $\epsilon,\nu_0>0$. Choose $J$ uniformly as above. The
$j>J$ tail is bounded independently of both imaginary parts by
\eqref{eq:logP}. Each of the finitely many factors with $j\le J$
has real part in a bounded inner strip of $\Rea w>\beta$.
The Hecke strip bound and Lemma~\ref{am:reciprocal-growth} bound their
product by a fixed power of
$1+|\Ima s|+|\Ima v|$, since $j\le J$.
The same lemma bounds $1/P(s)$, and $E_\infty$ is uniformly bounded.
This proves the polynomial-growth assertion for $\HH$.

For real $s\to+\infty$, the direct product \eqref{eq:H0} gives
\begin{equation}
 \sum_{p\text{ split}}|h_p(s,v)-1|
 \le C_{\nu_0}\sum_p p^{-s}\longrightarrow0
 \quad\text{uniformly for }\Rea v\ge\nu_0.
\end{equation}
Dominated convergence applies using $s\ge2$ as a summable
majorant. The elementary product estimate
\begin{equation}
 \left|\prod_p(1+u_p)-1\right|
 \le\exp\left(\sum_p|u_p|\right)-1
\end{equation}
then proves \eqref{eq:limit}, uniformly over the entire unbounded
half-plane in $v$. This completes Theorem~\ref{thm:mixed}.

\section{Poles that survive at one fixed norm parameter}

Suppose $A$ has a pole of order $k$ at $v_0$ with $\Rea v_0>0$.
Choose $0<\nu_0<\Rea v_0$. The uniform limit
\eqref{eq:limit} makes $|\HH(s,v)-1|<1/2$ throughout
$\Rea v\ge\nu_0$ for all sufficiently large real $s$.
In particular $\HH(s,v_0)\ne0$. Since $\HH$ is holomorphic
near $v_0$, multiplication by it preserves the pole order $k$.

For the stronger assertion at $s=4$, use only the direct
split-prime product \eqref{eq:H0}, which exists without an extra
zero-free hypothesis. Every split prime is at least $7$.
For $\Rea v\ge1/4$,
\begin{equation}\label{eq:fixedbound}
 |h_p(4,v)-1|
 \le\frac{2p^{-4}}{1-p^{-1/4}}
 \le8p^{-4}\le\frac{8}{7^4}<1.
\end{equation}
Here $7^{-1/4}<3/4$, since $7>256/81$, so
$1-p^{-1/4}>1/4$. Every local factor is therefore nonzero.
The bounds in \eqref{eq:fixedbound} are summable, uniformly in $v$,
and all factors lie in a fixed disk about $1$ of radius less than
$1$. Their logarithms converge absolutely and uniformly.
Their product is consequently nowhere zero on
$\Rea v\ge1/4$, and is holomorphic on $\Rea v>0$.
This proves \eqref{eq:s4} and preservation of the full order of
every pole there. No assertion about zeros of $P$ to the left of
its ordinary Euler half-plane was used.

\section{Counting every possible numerator cancellation}\label{sec:count}

Let $Z_R(2T)$ count zeta zeros $\rho$ with
$\Rea\rho\ge1/2$ and $0<\Ima\rho\le2T$, with multiplicity.
The functional equation and real coefficients map $\rho$ to
$1-\overline\rho$ with the same positive ordinate and multiplicity.
Zeros off the critical line are paired between the two halves;
critical-line zeros are retained in $Z_R$. Thus
\begin{equation}\label{eq:rightcount}
 Z_R(2T)\ge\frac12N_\zeta^+(2T).
\end{equation}
The substitution $\rho=2v$ identifies these zeros with denominator
zeros of \eqref{eq:A} in \eqref{eq:region}, preserving multiplicity.
Neither finite correction factor vanishes when $\Rea v>0$,
and the pole of $\zeta(2v)$ at $v=1/2$ lies outside that region.
If $m_\zeta(2v)$ and $m_\eta(v)$ denote the respective zero orders,
the pole order at such a point is
\begin{equation}\label{eq:multiplicity}
 \max\{m_\zeta(2v)-m_\eta(v),0\}.
\end{equation}
Summing and using $\max(a-b,0)\ge a-b$ gives
\begin{equation}\label{eq:subtract}
 Q_A(T)\ge Z_R(2T)-C_\eta(T),
\end{equation}
where $C_\eta(T)$ counts all numerator-zero multiplicities in
\eqref{eq:region}; including numerator zeros that do not coincide
with denominator zeros only weakens this inequality.

A zero $v$ counted in $C_\eta$ has $\Ima v>0$ and
$0<\Rea v<1/2$. The real primitive character $\eta$ has the
functional-equation and conjugation symmetries that produce
\begin{equation}\label{eq:quartet}
 v,\qquad\overline v,\qquad1-v,\qquad1-\overline v.
\end{equation}
These four zeros are distinct, have equal multiplicity, and lie
in $0<\Rea v<1$, $|\Ima v|\le T$. Exactly one member of such an
orbit can lie in \eqref{eq:region}; distinct counted zeros belong
to disjoint orbits. Zeros on the critical line or real axis can
only increase the total $N_\eta^{\pm}(T)$ without affecting this
argument. Therefore
\begin{equation}\label{eq:quarter}
 C_\eta(T)\le\frac14N_\eta^{\pm}(T).
\end{equation}
Combining \eqref{eq:rightcount}, \eqref{eq:subtract}, and
\eqref{eq:quarter} proves \eqref{eq:count}. The lower boundary
$\Rea v=1/4$ and the height endpoint $\Ima v=T$ are included.
The strict upper bound $\Rea v<1/2$ ensures that the numerator
quartet is distinct. Arbitrary common zeros and arbitrary
multiplicities have been allowed throughout.

\paragraph{The exact external estimates.}
Hasanalizade, Shen, and Wong (2022, Corollary~1.2, printed p.~221)
prove
\begin{equation}\label{eq:HSW}
 \left|N_\zeta^+(t)-\frac{t}{2\pi}
             \log\frac{t}{2\pi e}\right|
 \le E_\zeta(t)\qquad(t\ge e).
\end{equation}
Their definition on printed p.~220 counts positive ordinates.
Bennett, Martin, O'Bryant, and Rechnitzer (2021, Corollary~1.2,
printed p.~1456), at conductor $3$, prove
\begin{equation}\label{eq:BMOR}
 \left|N_\eta^{\pm}(t)-\frac{t}{\pi}
                  \log\frac{3t}{2\pi e}\right|
 \le E_\eta(t)\qquad(t\ge5/7).
\end{equation}
Their printed p.~1455 explicitly defines the count using both
signs of the ordinate and multiplicities. These two estimates
are imported as published theorems. The present proof derives
the symmetry, cancellation, and pole-count consequences; it does
not claim to reconstruct the analytic proofs of
\eqref{eq:HSW} or \eqref{eq:BMOR}.

If $T\ge e/2$, then $2T\ge e$ and $T\ge5/7$.
Use the lower estimate in \eqref{eq:HSW} at $2T$ and the upper
estimate in \eqref{eq:BMOR} at $T$. The main term simplifies to
\begin{align}
 \frac{T}{2\pi}\log\frac{T}{\pi e}
 -\frac{T}{4\pi}\log\frac{3T}{2\pi e}
 &=\frac{T}{4\pi}\log\frac{2T}{3\pi e}.
\end{align}
The errors enter with coefficients $1/2$ and $1/4$, exactly as
in \eqref{eq:explicit}. Both errors are $O(\log T)$, proving
\eqref{eq:asymp} and showing that $Q_A(T)$ tends to infinity.
This completes Theorem~\ref{thm:count}.

\section{The precise continuation barrier}

Theorem~\ref{thm:count} gives infinitely many poles of $A$ in
$1/4\le\Rea v<1/2$ with positive imaginary part after every
possible numerator cancellation. Theorem~\ref{thm:poles} makes
each of them a pole of $\F(4,v)$ of exactly the same order.
For $u_0<1/4$, all these points lie in $\Rea v>u_0$, while
$s=4$ lies in $\Rea s>\beta$ because $\beta\le1$.
If a holomorphic continuation to the whole proposed tube existed,
its restriction at $s=4$ would agree with $A(v)\HH_0(4,v)$ by the
identity theorem, starting in $\Rea v>1$. It could not be
holomorphic at the proved poles. This proves the unconditional
full-tube obstruction.

Under zeta RH, all nontrivial zeros of $\zeta(2v)$ in the positive
real half-plane lie on $\Rea v=1/4$. Its trivial zeros are on the
negative real axis. Its pole at $v=1/2$ is a zero of
$1/\zeta(2v)$ and causes no pole of $A$. The numerator is entire
and the finite correction factors have no zeros for $\Rea v>0$.
Thus $A$ is holomorphic on $\Rea v>1/4$. The cancellation-inclusive
count still proves infinitely many boundary poles, and every one
survives at $s=4$.

These arguments concern continuation holomorphic on an entire
tube. They neither estimate the residues nor establish that
residue contributions cannot cancel in another transform or
another sum. They give no stronger zero-free strip for $P$,
no cubic RH theorem, and no unconditional assertion of zeta RH.
No stronger bound for the cubic balanced prime sum is obtained. This completes Corollary~\ref{cor:barrier}.


\section{Smooth Mellin bounds and the original joint tail}

\begin{lemma}[Uniform smooth Mellin inversion]\label{res:smooth}
For fixed \(\sigma>0\), \(\sigma\leq\Rea s_i\leq2\), and every
fixed \(J\), one has
\[
 |\widehat\Psi_q(s_1,s_2)|
 \ll_J q^{-\sigma}\log(2+q)
 \prod_i(1+|\Ima s_i|)^{-J}.
\]
The transform \(\widetilde K\) has arbitrary inverse-power decay
on every fixed positive real-part line. The initial identity is
\begin{align}\label{res:initial-mellin}
 J_V=\frac1{(2\pi i)^3}
 \int_{(z)}\int_{(2)}\int_{(2)}
 &\widetilde K(z)R^{-z}U^{s_1+s_2}
 \widehat\Psi_q(s_1,s_2)\\
 &{}\cdot\F(s_1,2s_1-z)\F(s_2,2s_2-z)\,ds_2\,ds_1\,dz,
 \notag
\end{align}
where \(0<\Rea z<3\), and all initial interchanges are absolute.
\end{lemma}
\begin{proof}
On the support of \(\Psi_q\), both \(t_i\leq2\), their product
lies in a fixed compact subinterval of \((1/q,2/q)\), and each
\(t_i\gg1/q\). The logarithmic area of this region is
\(O(\log(2+q))\). If \(\sigma_i\in[\sigma,2]\), then
\[
 t_1^{\sigma_1}t_2^{\sigma_2}
 =(t_1t_2)^\sigma
 t_1^{\sigma_1-\sigma}t_2^{\sigma_2-\sigma}
 \ll q^{-\sigma}.
\]
In logarithmic coordinates, every fixed mixed derivative of the
weight has the same support and uniform boundedness: derivatives
of \(V(qt_1t_2)\) involve its logarithmic derivatives at a point
in the fixed support, and those of \(b(t_i)\) are bounded.
Repeated integration by parts in each coordinate therefore gives
the asserted rapid decay with only finitely many seminorms.
The boundary terms vanish, since the logarithmic support is compact
for each \(q\).
Also
\[
 \widetilde K(z)=\frac1z\int_1^2K'(t)t^{-z}\,dt,
\]
and the last integral is the Fourier transform of a fixed smooth
compactly supported function in logarithmic coordinates.
Repeated integration by parts proves its decay.
Three ordinary Mellin inversions now give
\eqref{res:initial-mellin}. On the indicated initial lines,
\(\Rea s_i=2\) and \(\Rea(2s_i-z)=4-\Rea z>1\).
The absolutely convergent ideal series from \eqref{eq:F},
together with the rapid transform bounds, justify summation and
integration in either order.

For \(V_n\), derivatives of the cutoff occur only when
\((L_X/n)t\in[1,2]\). On the fixed support of \(W\), that condition
bounds \(L_X/n\) above and below by positive constants.
Outside it the cutoff is constant. This proves the uniform finite
seminorm assertion used for the balanced application.
\end{proof}

\begin{lemma}[Joint-tail estimate]\label{res:oldtail}
Under the two zero-free hypotheses above,
\[
 J_V(Y,U,R)\ll_\varepsilon
 Y^{\beta+\varepsilon}R^{-(2\beta-u)}\log(2+q).
\]
\end{lemma}
\begin{proof}
Take small fixed \(\delta,\eta>0\), put
\(\sigma=\beta+\delta\), \(v_*=u+\eta<1/2\), and
\(a=2\sigma-v_*\in(0,3)\). Use \(\Rea z=a\) in
\eqref{res:initial-mellin} and move both norm contours from \(2\)
to \(\sigma\). Along these deformations
\(\Rea(2s_i-z)\geq v_*>u\).
Theorem~\ref{thm:mixed} and the zero-free hypothesis for zeta
therefore make every integrand holomorphic in the traversed region.
The reciprocal bound follows from
Lemma~\ref{am:reciprocal-growth} applied to
\((w-1)\zeta(w)\), and all remaining growth is polynomial.
Lemma~\ref{res:smooth} supplies an integrable majorant and makes
the horizontal edges vanish. On the final contours the real power is
\[
 U^{2\sigma}q^{-\sigma}R^{-a}
 =Y^\sigma R^{-(2\sigma-v_*)}.
\]
Since \(r(\mathfrak a)^2\leq\N\mathfrak a\), the summand is zero
for \(R\geq\sqrt{2Y}\): on the support of \(V\) the product of
contents is strictly less than \(\sqrt{2Y}\), where \(K(t)=0\)
for \(t\leq1\). In the remaining range \(1\leq R<\sqrt{2Y}\),
\[
 Y^\delta R^{-2\delta+\eta}\ll_\eta
 Y^{\delta+\eta/2}.
\]
Choose the fixed gaps so that \(\delta+\eta/2<\varepsilon\).
This proves the stated uniform estimate, including arbitrarily
large \(R\), where the tail vanishes.
\end{proof}


\section{Good heights and complete grouped
residues}\label{good-heights-and-complete-grouped-residues}

\subsection{Location and orders of the positive-real-part
poles}\label{location-and-orders-of-the-positive-real-part-poles}

For \(\Re v>0\), neither \(1-3^{-2v}\) nor \(1-269^{-v}\) vanishes,
since each corresponding exponential has modulus strictly less than one.
The numerator \(L(v,\eta)\) is entire. The pole of \(\zeta(2v)\) at
\(v=1/2\) makes \(1/\zeta(2v)\) vanish and does not produce a pole of
\(A\).

Consequently a pole of \(A\) with positive real part must come from a
nontrivial zero of \(\zeta(2v)\). The zero-free hypothesis gives
\(\Re(2v)\leq\rho\). The functional equation reflects nontrivial zeros
across \(\Re s=1/2\), with multiplicity. A zero with \(\Re(2v)<1-\rho\)
would therefore yield one with real part greater than \(\rho\), which is
forbidden. Thus

\[
 \ell\leq\Re v\leq u
\]

for every possible positive-real-part pole. At any candidate point,
zeros of \(L(v,\eta)\) cancel the corresponding part of the denominator
order. A canceled point is not a pole. All statements below refer to the
resulting meromorphic function \(A\), so these cancellations are
included automatically.

If a genuine pole \(v_0\) has order \(m\) and local principal part \[
 A(v)=\sum_{j=1}^{m}a_{-j}(v-v_0)^{-j}+O(1),
\] then, for holomorphic \(\Phi\),

\begin{equation}\label{res:A1}
\operatorname{Res}_{v=v_0}(A(v)\Phi(v))
 =\sum_{j=1}^{m}\frac{a_{-j}}{(j-1)!}\Phi^{(j-1)}(v_0).
\end{equation}

This formula includes multiple poles and further cancellation by
\(\Phi\). It is not estimated through individual inverse derivatives of
\(\zeta\).

\subsection{Polynomial bounds on the two vertical
lines}\label{polynomial-bounds-on-the-two-vertical-lines}

Put \(f(s)=(s-1)\zeta(s)\). This function is holomorphic and zero-free
on \(\Re s>\rho\), including \(s=1\). It has polynomial growth in fixed
strips. On \(\Re s=2\), the Euler-product lower bound for \(\zeta\)
supplies the lower bound needed for
Lemma~\ref{am:reciprocal-growth}. That lemma therefore
gives polynomial bounds for \(1/f\), and hence for
\(1/\zeta(s)=(s-1)/f(s)\), on every closed inner real-part strip in
\(\Re s>\rho\).

Since \(2v_R>\rho\), this gives \[
 \zeta(2v_R+2it)^{-1}\ll (1+|t|)^C.
\] For the left line, \(0<2v_L<1-\rho\). The functional equation writes
\[
 \zeta(s)=\chi(s)\zeta(1-s).
\] On the fixed line \(\Re s=2v_L\), the gamma and trigonometric factor
\(\chi\) is nonzero; its modulus and reciprocal have polynomial growth
for large imaginary part, and are bounded on the remaining compact part.
The reflected line has real part \(1-2v_L>\rho\). The preceding
reciprocal bound therefore also gives \[
 \zeta(2v_L+2it)^{-1}\ll (1+|t|)^C.
\]

The numerator of \(A\) has polynomial growth by the stated analytic
input. The two elementary denominator factors are bounded away from zero
uniformly for \(v_L\leq\Re v\leq v_R\). Hence

\begin{equation}\label{res:A2}
|A(v_L+it)|+|A(v_R+it)|
 \ll (1+|t|)^C.
\end{equation}

The same argument gives a polynomial bound uniformly on any fixed strip
\(v_R\leq\Re v\leq w\). This latter observation is used when moving the
initial right contour.

\subsection{A good-height construction that does not require zero
spacing}\label{a-good-height-construction-that-does-not-require-zero-spacing}

Here is a direct minimum-modulus argument with an averaged logarithm. It
avoids the inadequate inference that merely staying a
reciprocal-logarithmic distance from all zeros must give a polynomial
reciprocal bound.

Set \(a=2v_L\), \(b=2v_R\), so \(0<a<b<1\). For a sufficiently large
integer \(N\), center a disk at \[
 c_N=2+i(2N+1).
\] On \(|s-c_N|\leq4\), the zeta function is holomorphic, and the
fixed-strip growth input gives \[
 |\zeta(s)|\leq M_N,\qquad
 M_N\geq 1,\qquad \log M_N\ll\log(2+N).
\] The Euler-product lower bound gives \(|\zeta(c_N)|\geq c>0\).

Jensen's formula, or its immediate disk inequality, shows that the
number of zeros in \(|s-c_N|<13/4\), counted with multiplicity, is
\(O(\log(2+N))\). Indeed each zero in that smaller disk contributes at
least \(\log(4/(13/4))>0\) to the Jensen sum, whose total is at most
\(\log M_N-\log|\zeta(c_N)|\).

Choose a radius \(r_N\in[3,13/4]\) on whose circle there is no zero. Let
\(\rho_j=c_N+a_j\) be all zeros inside this circle, repeated with
multiplicity, and form the finite disk Blaschke product \[
 B_N(w)=\prod_j \frac{r_N(w-a_j)}{r_N^2-\overline{a_j}w},
\] with harmless unimodular factors if desired. It has modulus one on
\(|w|=r_N\). The function \[
 g_N(w)=\frac{\zeta(c_N+w)}{B_N(w)}
\] is holomorphic and zero-free in that disk. The maximum principle
gives \(|g_N|\leq M_N\), while \[
 |g_N(0)|=\frac{|\zeta(c_N)|}{\prod_j |a_j|/r_N}
 \geq|\zeta(c_N)|\geq c.
\]

The function \(h_N=\log M_N-\log|g_N|\) is nonnegative and harmonic.
Harnack's inequality on \(|w|\leq5/2\), with the outer radii uniformly
at least \(3\), gives \(h_N(w)\ll h_N(0)\ll\log(2+N)\). Consequently

\begin{equation}\label{res:A3}
|g_N(w)|\geq(2+N)^{-C}
 \quad (|w|\leq5/2).
\end{equation}

For \(x\in[a,b]\) and \(t\in I_N=[N+1/4,N+3/4]\), the point \(s=x+2it\)
lies in \(|s-c_N|<5/2\). On this smaller disk, \[
 \left|\frac{r_N(w-a_j)}{r_N^2-\overline{a_j}w}\right|
 \geq \frac{|s-\rho_j|}{6}.
\] Moreover, \(|s-\rho_j|\geq|2t-\Im\rho_j|\). For every real \(y\),

\begin{equation}\label{res:A4}
\int_{I_N}\log^+\frac{6}{|2t-y|}\,dt
 \leq \frac12\int_{-6}^{6}\log\frac{6}{|x|}\,dx=6.
\end{equation}

Summing \eqref{res:A4} over the \(O(\log(2+N))\) zeros, there is a
\(T_N\in I_N\), away from their ordinates, for which \[
 \sum_j\log^+\frac{6}{|2T_N-\Im\rho_j|}
 \ll\log(2+N).
\] This bound is independent of \(x\in[a,b]\). Combining it with
\eqref{res:A3} proves

\begin{equation}\label{res:A5}
\sup_{a\leq x\leq b}|\zeta(x+2iT_N)^{-1}|
 \ll(2+N)^C.
\end{equation}

Small nonnegative \(N\) can be handled by choosing a height in the
specified interval outside the finite set of relevant zero ordinates and
enlarging the constant. Complex conjugation gives the same bounds at
\(-T_N\). Multiplying by the polynomially bounded numerator and the
uniformly harmless elementary factors proves

\begin{equation}\label{res:A6}
\sup_{v_L\leq x\leq v_R}
 \bigl(|A(x+iT_N)|+|A(x-iT_N)|\bigr)
 \ll(2+N)^C.
\end{equation}

The argument only used a local zero count and the integrability of a
logarithmic singularity. It is valid for arbitrarily close zeros and
arbitrary multiplicities.

\subsection{The block functionals}\label{the-block-functionals}

Define \(Q_0\) to have vertical sides \(v_L,v_R\) and horizontal sides
\(-T_0,T_0\). For \(N\geq1\), let \(Q_N\) have horizontal sides
\(T_{N-1},T_N\), and let \(Q_{-N}\) be its reflection in the real axis.
Each rectangle has bounded perimeter, and every point of \(Q_N\) has
imaginary part \(N+O(1)\). Their interiors are disjoint, their union
fills the strip, and no pole lies on any boundary.

For \(\Phi\) holomorphic near \(\overline{Q_N}\), define

\begin{equation}\label{res:A7}
\mathcal R_N[\Phi]
 =\sum_{v_0\in Q_N}\operatorname{Res}_{v=v_0}(A(v)\Phi(v))
 =\frac{1}{2\pi i}\int_{\partial Q_N}^{+}A(v)\Phi(v)\,dv.
\end{equation}

The sum in an individual rectangle is finite. The superscript \(+\)
means positive orientation. By \eqref{res:A2}, \eqref{res:A6}, and
bounded perimeter, there are fixed constants \(C_0,C_1\) such that

\begin{equation}\label{res:A8}
|\mathcal R_N[\Phi]|
 \leq C_0(1+|N|)^{C_1}
 \sup_{\partial Q_N}|\Phi|.
\end{equation}

This is a bound for the complete residue functional of a rectangle. It
includes every multiplicity and every numerator or test-function
cancellation. It does not bound the sum of the absolute values of its
individual residues. All infinite residue sums below are sums of the
complex numbers \eqref{res:A7}, and absolute convergence means
\(\sum_N|\mathcal R_N[\Phi]|<\infty\). No absolute convergence of the
individual zero residues is asserted.

\section{Exact two-variable decomposition, absolute block convergence,
and
powers}\label{exact-two-variable-decomposition-absolute-block-convergence-and-powers}

\subsection{A common integrable
majorant}\label{a-common-integrable-majorant}

Let \(\tau=\Im z\), \(t_j=\Im v_j\). On any of the fixed real-part
strips used below, \[
 \Im((z+v_j)/2)=(\tau+t_j)/2.
\] The polynomial growth of \(\HH\), the stated Mellin bounds, and
arbitrary inverse-power decay of \(\widetilde K\) imply that, for any
desired large \(J\), the absolute value of the kernel is bounded by its
real-power factor times

\begin{equation}\label{res:B1}
C_J(1+|\tau|+|t_1|+|t_2|)^C
 (1+|\tau|)^{-J}
 (1+|\tau+t_1|)^{-J}
 (1+|\tau+t_2|)^{-J}.
\end{equation}

Polynomial bounds for \(A\) on the relevant contours can be absorbed by
increasing \(C\).

The linear map \[
 (\tau,t_1,t_2)\longmapsto(\tau,\tau+t_1,\tau+t_2)
\] is invertible. Thus, after increasing \(J\), \eqref{res:B1} is
integrable over all three real variables. The corresponding majorant
remains summable if one or both \(t_j\) are replaced by a block index
\(N\), because a point on \(\partial Q_N\) has imaginary part
\(N+O(1)\). More explicitly, after absorbing the polynomial factor,

\begin{equation}\label{res:B2}
\int_{\mathbb R}(1+|\tau|)^{-p}
 \sum_N(1+|\tau+N|)^{-p}
 \sum_M(1+|\tau+M|)^{-p}\,d\tau<\infty
 \quad(p>1),
\end{equation}

since each lattice sum is bounded uniformly in \(\tau\). Replacing
either sum by an integral gives the single-block and no-block versions.

For horizontal-edge limits, the product \[
 (1+|\tau|)^{-p}(1+|\tau+T|)^{-p}
\] has an integral tending to zero as \(|T|\to\infty\), even after any
fixed polynomial factor has been absorbed into a larger choice of \(p\).
This follows by splitting into \(|\tau|\geq |T|/2\) and its complement.
Consequently the horizontal edges also vanish after integrating the
other free variables. These facts justify the contour limits and every
Fubini or Tonelli interchange below.

Only finitely many derivative seminorms are needed for each fixed
polynomial growth exponent. The strong derivative condition on \(W\) is
compatible with those requirements; no extra power saving is inferred
from that condition.

\subsection{Moving to the common right
line}\label{moving-to-the-common-right-line}

In the initial Mellin inversion in
Lemma~\ref{res:smooth}, choose \(\Re z=\alpha\) at the
outset. This is allowed because \(0<\alpha<3\). For fixed \(z\), put \[
 v_j=2s_j-z,\qquad ds_j=dv_j/2.
\] The initial real parts are \[
 \Re v_j=w_0:=4-\alpha>1,\qquad \Re s_j=2.
\] The two Jacobians give the factor \(1/4\). Therefore

\begin{equation}\label{res:B3}
J_V=\frac1{4(2\pi i)^3}
 \int_{(\alpha)}\int_{(w_0)}\int_{(w_0)}
 A(v_1)A(v_2)\mathcal K_z(v_1,v_2)\,dv_2\,dv_1\,dz.
\end{equation}

Between \(v_R\) and \(w_0\), \(A\) is holomorphic and polynomially
bounded. The kernel is holomorphic in each \(v_j\), since its
\(s_j\)-real part is at least \(\sigma>\beta\) and its \(v_j\)-real part
is positive. In fact every \(s_j\)-real part encountered lies between
\(\sigma\) and \(2\), as required for the stated Mellin bound. The
majorant in the common-majorant argument permits moving both
\(v_j\)-lines in \eqref{res:B3} from \(w_0\) to \(v_R\), with no
residue. There has been no shift of the \(z\)-line.

\subsection{The exact four-term
identity}\label{the-exact-four-term-identity}

For fixed \(z\), shift one \(v_j\)-line from \(v_R\) to \(v_L\) through
finite rectangles with the good horizontal edges. Positive orientation
gives

\begin{equation}\label{res:B4}
\int_{(v_R)} A(v)\Phi(v)\,dv
 =
 \int_{(v_L)} A(v)\Phi(v)\,dv
 +2\pi i\sum_N\mathcal R_N[\Phi].
\end{equation}

The limiting identity follows from the edge estimates and absolute
convergence of the complete blocks.

Apply \eqref{res:B4} in both variables of \eqref{res:B3}. This gives
exactly

\begin{equation}\label{res:B5}
J_V=I_L+S_1+S_2+D,
\end{equation}

where

\begin{equation}\label{res:B6}
I_L=\frac1{4(2\pi i)^3}
 \int_{(\alpha)}\int_{(v_L)}\int_{(v_L)}
 A(v_1)A(v_2)\mathcal K_z(v_1,v_2)\,dv_2\,dv_1\,dz,
\end{equation}

\begin{equation}\label{res:B7}
S_1=\frac1{4(2\pi i)^2}
 \int_{(\alpha)}\int_{(v_L)}
 A(v_2)\sum_N\mathcal R_N^{(1)}
       [\mathcal K_z(\,\cdot\,,v_2)]\,dv_2\,dz,
\end{equation}

\(S_2\) is the symmetric expression, and

\begin{equation}\label{res:B8}
D=\frac1{4(2\pi i)}
 \int_{(\alpha)}
 \sum_{N,M}\mathcal R_N^{(1)}\mathcal R_M^{(2)}
       [\mathcal K_z]\,dz.
\end{equation}

For the double functional, its two compact contour integrals commute by
ordinary finite-contour Fubini:

\begin{equation}\label{res:B9}
\mathcal R_N^{(1)}\mathcal R_M^{(2)}[\mathcal K_z]
 =
 \frac1{(2\pi i)^2}
 \int_{\partial Q_N}^{+}\int_{\partial Q_M}^{+}
 A(v_1)A(v_2)\mathcal K_z(v_1,v_2)\,dv_2\,dv_1.
\end{equation}

Applying \eqref{res:A8} twice, then \eqref{res:B1}--\eqref{res:B2},
proves

\begin{equation}\label{res:B10}
\int_{(\alpha)}
 \sum_{N,M}
 \left|\mathcal R_N^{(1)}\mathcal R_M^{(2)}
       [\mathcal K_z]\right|\,|dz|<\infty,
\end{equation}

with their fixed real-power factors understood. The analogous absolute
estimates hold for both singles and the bulk. Thus the order of the two
block sums, their order relative to the remaining contour integrals, and
the two contour-shift orders can all be interchanged. This is stronger
than conditional convergence of the displayed grouped sums, but it still
makes no claim of termwise absolute convergence over individual zeros.

\subsection{Bounds, including the full real-power
costs}\label{bounds-including-the-full-real-power-costs}

Write \(\Delta=v_R-v_L>0\). Lemma~\ref{res:smooth} is
used with its lower real-part parameter \(\sigma\): \[
 |\widehat\Psi_q(s_1,s_2)|
 \ll q^{-\sigma}\log(2+q)
 \prod_j(1+|\Im s_j|)^{-J}.
\]

For the bulk, \(\Re v_1=\Re v_2=v_L\), so \[
 U^{\Re(z+(v_1+v_2)/2)}q^{-\sigma}
 =U^{\alpha+v_L}q^{-\sigma}
 =U^{2\sigma}q^{-\sigma}=Y^\sigma.
\] All imaginary-variable integrals are bounded by the common-majorant
argument. Therefore

\begin{equation}\label{res:B11}
I_L\ll Y^\sigma R^{-\alpha}\log(2+q).
\end{equation}

For a single residue, one variable ranges over a block boundary with
real part at most \(v_R\), and the other is on \(v_L\). The
corresponding power is at most \[
 U^{2\sigma+\Delta/2}q^{-\sigma}
 =Y^{\sigma+\Delta/4}q^{\Delta/4}.
\] Hence

\begin{equation}\label{res:B12}
|S_1|+|S_2|
 \ll Y^{\sigma+\Delta/4}q^{\Delta/4}
 R^{-\alpha}\log(2+q).
\end{equation}

For the double residues, both variables range up to \(v_R\). The power
is at most \[
 U^{2\sigma+\Delta}q^{-\sigma}
 =Y^{\sigma+\Delta/2}q^{\Delta/2}.
\] Consequently

\begin{equation}\label{res:B13}
D\ll Y^{\sigma+\Delta/2}q^{\Delta/2}
 R^{-\alpha}\log(2+q).
\end{equation}

The exponents \(\Delta/4\) and \(\Delta/2\) are part of the cost of
using complete block contours. They cannot be discarded on the ground
that residues have been taken.

\subsection{Balanced exponents and the meaning of a limiting
exponent}\label{balanced-exponents-and-the-meaning-of-a-limiting-exponent}

In the balanced application, \[
 U^2=2X,\quad Y=X/n,\quad q=2n,\quad n<2L_X=X^{o(1)}.
\] Thus \(Y=X^{1-o(1)}\) and \(q=Y^{o(1)}\). With \(R=Y^\gamma\), first
choose fixed small \(\delta>0\) and \(v_R-u>0\), then absorb the
resulting arbitrarily small losses into \(Y^\varepsilon\). One must not
move contours at a gap depending on \(Y\) without controlling the
corresponding constants.

Taking the infima of the fixed-gap exponents gives

\begin{equation}\label{res:B14}
E_I=\beta-(2\beta-v_L)\gamma,
\end{equation}

\begin{equation}\label{res:B15}
E_S=\beta+\frac{u-v_L}{4}-(2\beta-v_L)\gamma,
\end{equation}

\begin{equation}\label{res:B16}
E_D^{\rm direct}
 =\beta+\frac{u-v_L}{2}-(2\beta-v_L)\gamma.
\end{equation}

In precise asymptotic language, each expression is an exponent valid
with an arbitrarily small positive \(\varepsilon\) loss.

\subsection{What the joint-tail bound implies for the complete double
functional}\label{what-the-joint-tail-bound-implies-for-the-complete-double-functional}

The exact identity gives

\begin{equation}\label{res:B17}
|D|\leq |J_V|+|I_L|+|S_1|+|S_2|.
\end{equation}

For arbitrary \(q\geq1\), the fully justified estimate is

\begin{equation}\label{res:B18}
|D|\ll\log(2+q)\left[
 Y^{\beta+\varepsilon}R^{-(2\beta-u)}
 +Y^\sigma R^{-\alpha}\{1+(Yq)^{\Delta/4}\}\right].
\end{equation}

This is the statement before using the balanced hypothesis.

If \(R\geq Y^{1/4}\), put \(\theta=u-v_L>0\), \(\eta=v_R-u>0\). Since
\(\alpha=(2\beta-u)+(2\delta+\theta)\),

\begin{equation}\label{res:B19}
Y^\sigma R^{-\alpha}\{1+(Yq)^{\Delta/4}\}
 \ll
 Y^{\beta+\delta/2+\eta/4}q^{\Delta/4}
 R^{-(2\beta-u)}
 \left(\frac{Y^{1/4}}R\right)^{2\delta+\theta}.
\end{equation}

The last factor is at most one. In the balanced regime \(q=Y^{o(1)}\),
selecting fixed small \(\delta,\eta\) yields, for every
\(\varepsilon>0\),

\begin{equation}\label{res:B20}
\boxed{\quad
 D\ll_\varepsilon
 Y^{\beta+\varepsilon}R^{-(2\beta-u)}\log(2+q),
 \qquad R\geq Y^{1/4}.
 \quad}
\end{equation}

This is a bound on the complete double-residue functional, with all
blocks included. It is not a bound of this strength on each individual
residue or on the sum of absolute values of individual residues.

Without \(q=Y^{o(1)}\), the factor in \eqref{res:B19} must be retained.
For example, one can state a general weakened version with
\(Y^\varepsilon q^{(u-v_L)/4+\varepsilon}\) multiplying the old scale.
The condition \(R\geq Y^{1/4}\) alone does not remove an unrestricted
\(q\)-power from the estimates supplied in the hypotheses above. This is
the precise qualification outside the balanced application.

\subsection{Why these estimates do not prove a smaller content
cutoff}\label{why-these-estimates-do-not-prove-a-smaller-content-cutoff}

Let

\begin{equation}\label{res:B21}
E_J(\gamma)=\beta-(2\beta-u)\gamma,\qquad
 \gamma_{\rm old}=\frac{\beta-1/2}{2\beta-u}.
\end{equation}

The denominator is positive and \(0<\gamma_{\rm old}<1/2\). For fixed
\(v_L\),

\begin{equation}\label{res:B22}
E_S-E_J=(u-v_L)(1/4-\gamma),\qquad
 E_D^{\rm direct}-E_J=(u-v_L)(1/2-\gamma).
\end{equation}

Thus, for \(1/4\leq\gamma\leq1/2\), the best double-functional exponent
furnished by \eqref{res:B13} and \eqref{res:B17} is \(E_J\), and
\eqref{res:B20} realizes it. For \(0\leq\gamma<1/4\), the triangle
inequality furnishes \(\max(E_J,E_S)=E_S\), still better than the direct
double bound but at least \(E_J\). Equivalently, on
\(0\leq\gamma\leq1/2\) the best exponent supplied by these two routes is

\begin{equation}\label{res:B23}
E_D^{\rm available}
 =E_J+(u-v_L)\max(0,1/4-\gamma).
\end{equation}

Whenever \(\gamma<\gamma_{\rm old}\), \(E_J>1/2\), so none of these
double-functional estimates reaches the square-root target. This remains
true on optimizing \(v_L\) within its permitted interval, because
\(u-v_L>0\).

This conclusion is a limitation of the displayed upper bounds. The
inequality \(E_D^{\rm available}>1/2\) is not a lower bound for \(|D|\),
\(|J_V|\), or any arithmetic sum. Cancellation between different blocks
or different cofactors could improve the actual complete functional.
Such an improvement is additional mathematics and has not been proved
here.

\section{Ordinary cofactors, direct cutoffs, and cap
cancellations}\label{ordinary-cofactors-direct-cutoffs-and-cap-cancellations}

\subsection{The full ordinary-cofactor
factor}\label{the-full-ordinary-cofactor-factor}

Let \[
 a(\mathfrak c)=-\frac12\psi(\mathfrak c)
 d_{2,K}(\mathfrak c)\log N\mathfrak c.
\] Since no cross-coprimality is imposed, multiplication of the
absolutely convergent supported Dirichlet series gives \[
 P(s)^2=\sum_{\mathfrak c}
 \psi(\mathfrak c)d_{2,K}(\mathfrak c)(N\mathfrak c)^{-s}.
\] Differentiating in its absolute half-plane yields

\begin{equation}\label{res:C1}
\sum_{\mathfrak c}a(\mathfrak c)(N\mathfrak c)^{-s}
 =P(s)P'(s).
\end{equation}

The factor \(1/2\) and the sign are both accounted for.

With a global content cutoff independent of \(N\mathfrak c\), and with
no \(n\)-dependent direct weight, \eqref{res:C1} is the full ordinary
cofactor. For an uncapped pair, putting \(v=2s-z\) gives

\begin{equation}\label{res:C2}
P(s)P'(s)\F(s,2s-z)^2
 =A(2s-z)^2G(s,2s-z)^2\,\frac{P'(s)}{P(s)}.
\end{equation}

Thus the double denominator \(P(s)^2\) is reduced to a logarithmic
derivative, not automatically removed.

For clarity, the domain of Theorem~\ref{thm:mixed}
itself contains no zero of \(P\). If a local continuation to a point
\(s_0\) where \(P\) has order \(m\geq1\) is available and the remaining
composed factors are holomorphic there, then \[
 \frac{P'(s)}{P(s)}=\frac{m}{s-s_0}+O(1).
\] The pole of this quotient is simple even for a multiple zero of
\(P\). If \(G(s,2s-z)\) has order \(g\geq0\) at that point and
\(A(2s-z)\) is regular and nonzero, the total order in \eqref{res:C2} is
\(2g-1\). A zero of \(G\) can therefore remove the
logarithmic-derivative pole. Zeros or poles of \(A\) can modify the
total order as well. Nothing in the hypotheses above proves \(g=0\), so
a genuine surviving pole cannot be asserted unconditionally. Conversely,
nothing proves that the zeros of \(G\) cancel all such factors.

For the actual capped pair, the factor is \(P(s)P'(s)M_U(s,z)^2\), not
\eqref{res:C2}. The cap expansion must be kept before inferring analytic
or power consequences.

\subsection{The cofactor-dependent content
cutoff}\label{the-cofactor-dependent-content-cutoff}

If \[
 R_{\mathfrak c}=(X/N\mathfrak c)^\gamma,
\] then the two Mellin weights give \[
 (X/n)^sR_{\mathfrak c}^{-z}
 =X^{s-\gamma z}n^{-(s-\gamma z)}.
\] Consequently the full untruncated cofactor is

\begin{equation}\label{res:C3}
P(s-\gamma z)P'(s-\gamma z).
\end{equation}

It does not in general cancel the denominator \(P(s)^2\) in an uncapped
pair. Equality of the two arguments, or coincident zeros, would be a
special condition rather than a general identity.

\subsection{What the direct cutoff actually
does}\label{what-the-direct-cutoff-actually-does}

In the balanced sum, \[
 V_n(t)=W(t)\left[1-\chi_d((L_X/n)t)\right].
\] Because \(W\) is supported in \((1,2)\), this vanishes identically
for \(n\geq2L_X\). Let \[
 \widehat V_n(s)=\int_0^\infty V_n(t)t^s\,\frac{dt}{t}.
\] The exact cofactor expression with the direct cutoff is

\begin{equation}\label{res:C4}
\mathcal C_{L_X}(s,z)
 =\sum_{N\mathfrak c<2L_X}
 a(\mathfrak c)(N\mathfrak c)^{-(s-\gamma z)}
 \widehat V_{N\mathfrak c}(s).
\end{equation}

It is a finite sum of entire functions of \(s,z\). The weighted capped
tail has the exact two-Mellin representation

\begin{equation}\label{res:C5}
\mathscr J(X)=\frac1{(2\pi i)^2}
 \int_{(z)}\int_{(s)}
 \widetilde K(z)X^{s-\gamma z}
 \mathcal C_{L_X}(s,z)M_U(s,z)^2\,ds\,dz.
\end{equation}

Here the initial lines can be taken in an absolute domain; finite caps
and a finite cofactor sum make the displayed expression unproblematic
there. Formula \eqref{res:C4}, rather than a full Euler product, is the
applicable cofactor in this form of the balanced calculation.

There is a second useful exact description. Put \(K_d(t)=1-\chi_d(t)\).
If \(A_0=N\mathfrak a_1N\mathfrak a_2\), then \[
 V_n(nA_0/X)=W(nA_0/X)K_d(L_XA_0/X).
\] The direct cutoff is independent of \(n\) in these original
variables, but it acts on the pair product. Mellin-inverting it with a
variable \(h\), \(\Re h>0\), gives an integrand proportional to

\begin{equation}\label{res:C6}
X^{s-h-\gamma z}L_X^h\,
 \widehat W(s)\widetilde K_d(h)\widetilde K(z)\,
 P(s-\gamma z)P'(s-\gamma z)\,
 M_U(s-h,z)^2.
\end{equation}

All inversions and the cofactor sum commute on sufficiently far-right
initial lines. In an uncapped version, the pair denominator in
\eqref{res:C6} is \(P(s-h)^2\). Hence the clean match in \eqref{res:C2}
does not persist term by term after adding this extra direct-cutoff
Mellin variable, even for \(\gamma=0\). Representations \eqref{res:C4}
and \eqref{res:C6} are consistent; they are two ways to retain the
cutoff that a naive use of \(P(s)P'(s)\) would otherwise omit.

\subsection{Meromorphic continuation of the cap
transform}\label{meromorphic-continuation-of-the-cap-transform}

For \(\Re w>0\),

\begin{equation}\label{res:C7}
\widehat b(w)=\int_0^\infty b(t)t^w\,\frac{dt}{t}
 =\frac1w+\int_1^2 b(t)t^w\,\frac{dt}{t}.
\end{equation}

The second term is entire, so \(\widehat b\) has precisely one pole, a
simple pole at \(w=0\) of residue one. Integration by parts also gives

\begin{equation}\label{res:C8}
\widehat b(w)=-\frac1w\int_1^2 b'(t)t^w\,dt.
\end{equation}

In logarithmic coordinates the last integrand is a smooth function
supported in a compact interval, with all endpoint derivatives
vanishing. Repeated integration by parts shows arbitrary inverse-power
decay on fixed vertical strips away from a neighborhood of \(w=0\).

Mellin inversion for \(b\), followed by absolute summation in a
sufficiently far-right domain, gives

\begin{equation}\label{res:C9}
M_U(s,z)=\frac1{2\pi i}\int_{(c)}
 \widehat b(w)U^w\F(s+w,2(s+w)-z)\,dw,
\end{equation}

where \(c>0\) is chosen so that \(\Re(s+w)>1\) and \(\Re(2(s+w)-z)>1\).
This is the cap-Mellin representation.

Independently of \eqref{res:C9}, the original definition of \(M_U\) is a
finite sum, since \(N\mathfrak d\leq2U\) wherever
\(b(N\mathfrak d/U)\ne0\). Thus \(M_U\) is entire in the two variables
\(s,z\). A pole in one term of a contour expansion of \eqref{res:C9}
must be canceled by the other terms.

\subsection{Moving content poles and complete cancellation at the cap
pole}\label{moving-content-poles-and-complete-cancellation-at-the-cap-pole}

Fix \(z\) locally and a content pole \(v_0\). Put \[
 \xi_0=(z+v_0)/2,\qquad
 f(\xi,z)=\F(\xi,2\xi-z).
\] In the permitted inner domain, \(\HH\) is holomorphic. Thus the
Laurent expansion at \(\xi_0\) has the form

\begin{equation}\label{res:C10}
f(\xi,z)=\sum_{j=1}^{m}c_{-j}(z)(\xi-\xi_0)^{-j}
          +f_0(\xi,z).
\end{equation}

Here \(m\) can be smaller than the order of the pole of \(A\), or zero,
because \(\HH\) may cancel it. A common maximum \(m\), allowing some
coefficients to vanish, works in a small neighborhood in \(z\).

In the \(w\)-integral the pole is at

\begin{equation}\label{res:C11}
w_0=\xi_0-s=(z+v_0)/2-s.
\end{equation}

It is moving because its position depends on \(s,z\). Away from
\(w_0=0\), its complete residue is

\begin{equation}\label{res:C12}
C_{v_0}(s,z;U)
 =\sum_{j=1}^{m}\frac{c_{-j}(z)}{(j-1)!}
 \left(\frac{d}{dw}\right)^{j-1}
 \bigl(\widehat b(w)U^w\bigr)\bigg|_{w=w_0}.
\end{equation}

The residue at \(w=0\) is \(f(s,z)\).

Write \[
 B_U(w)=\widehat b(w)U^w=\frac1w+h_U(w),
\] where \(h_U\) is entire, and set \(a=s-\xi_0\), so \(w_0=-a\). For
every \(j\geq1\),

\begin{equation}\label{res:C13}
\frac1{(j-1)!}
 \left(\frac{d}{dw}\right)^{j-1}\frac1w\bigg|_{w=-a}
 =(-1)^{j-1}(-a)^{-j}=-a^{-j}.
\end{equation}

Substituting into \eqref{res:C12} gives

\begin{equation}\label{res:C14}
C_{v_0}(s,z;U)
 =-\sum_{j=1}^{m}c_{-j}(z)a^{-j}
 +\sum_{j=1}^{m}\frac{c_{-j}(z)}{(j-1)!}
 h_U^{(j-1)}(-a).
\end{equation}

The second sum is holomorphic. The first is exactly the negative of the
full principal part of \(f(s,z)\). Therefore \[
 f(s,z)+C_{v_0}(s,z;U)
\] is holomorphic at the collision \(s=\xi_0\), for arbitrary pole
order. This cancels every negative Laurent coefficient, not just the
leading term.

At exact collision one can instead enclose both poles in one small
contour. Their combined residue is then the residue of the coincident
higher-order pole, and equals the holomorphic continuation of their sum.
Keeping both residues as a single local group is the clean
interpretation.

\subsection{Global cap block convergence and double-cap
combinations}\label{global-cap-block-convergence-and-double-cap-combinations}

The preceding local algebra is compatible with the infinite block
construction. On \(\Re z=\alpha\), content residues can be written in
the \(v\)-coordinate using \[
 w=(z+v)/2-s,\qquad dw=dv/2.
\] Away from a cap-pole collision, the content contribution of one block
is

\begin{equation}\label{res:C15}
\frac12\mathcal R_N\left[
 \widehat b((z+v)/2-s)\,
 U^{(z+v)/2-s}\HH((z+v)/2,v)\right].
\end{equation}

For fixed \(s,z,U\), or uniformly on a suitable compact set, the factor
\(\widehat b((z+v)/2-s)\) decays faster than any power of the block
height; \(\HH\) and the block norm grow polynomially. Thus the sum of
the absolute values of these complete block contributions converges.

There is one qualification to \eqref{res:C15}: if the transformed cap
pole \(v=2s-z\) lies in a rectangle, the bracket is not holomorphic
there. Its boundary integral includes the cap residue in addition to the
content residues. Subtract that cap residue when defining content-only
blocks, or group it with the content residues in that rectangle. Only
finitely many local blocks are affected, and at a collision the combined
group is the holomorphic expression proved in the cap-collision
calculation above. Thus one must not apply \eqref{res:A7} as a
holomorphic-test-function formula to a bracket that has an
unacknowledged pole.

To justify a cap contour shift directly, choose the left \(w\)-line so
that the corresponding \(v\)-real part lies strictly between \(0\) and
\(\ell\) and the \(\xi=s+w\)-real part remains greater than \(\beta\).
Around a collision this can be done because
\(\Re \xi_0=(\alpha+\Re v_0)/2>\sigma>\beta\). Choose the right line in
the initial absolute domain. The horizontal edges can be placed at the
transformed good heights \[
 \Im w=(\pm T_N+\Im z)/2-\Im s.
\] Polynomial growth and the rapid decay of \(\widehat b\) make those
edges vanish. The resulting left integral is locally holomorphic by
dominated convergence. The infinite content contributions converge in
complete blocks as just shown.

For two caps, the same arguments apply to the product of the two compact
contour integrals. The double block sum is absolutely convergent as a
sum of whole blocks. Compact-parameter versions give local holomorphy
after grouping collisions.

For an explicit justification with outer Mellin integrations, first keep
the common outer line at a fixed real part \(s_0\) so far right that
\(2s_0-\alpha>v_R+\kappa\), for some \(\kappa>0\), and the ordinary
cofactor series is absolute. Then \(\Re w_0\leq-\kappa/2\) on every
content block. The cap pole is outside the content strip, and all test
functions in \eqref{res:C15} are holomorphic there. Put
\(\omega=\Im s\). The rapid factors are \[
 (1+|\omega|)^{-J}(1+|\tau|)^{-J}
 \prod_{i=1}^2(1+|(\tau+t_i)/2-\omega|)^{-J},
\] from \(\widehat W\), \(\widetilde K\), and the two cap transforms.
The map from \((\omega,\tau,t_1,t_2)\) to the four arguments in these
factors is invertible. Polynomial growth is absorbed by increasing
\(J\), and the continuous-integral and lattice-sum argument of
\eqref{res:B2} proves the required joint absolute convergence. If the
extra direct-cutoff variable \(h\) is used, \(\widetilde K_d(h)\) adds
rapid decay in \(\Im h\), and the cap arguments use \(\omega-\Im h\);
the same invertible-coordinate argument applies. Thus both cap shifts,
the complete block sums, and the smooth outer integrations commute on
these specified contours. A later outer contour crossing a collision
must retain the combined cap/content group of the cap-collision
calculation above; this proof makes no claim that the separated singular
terms can be integrated across a collision without that grouping.

For the local cancellation itself, let \(f_i\) be the uncapped term of
cap \(i\), and \(C_i\) its matching moving residue. The four choices
generated by crossing both \(w=0\) poles and both corresponding moving
poles are

\begin{equation}\label{res:C16}
f_1f_2+f_1C_2+C_1f_2+C_1C_2
 =(f_1+C_1)(f_2+C_2).
\end{equation}

Each factor on the right is holomorphic at its collision by the
cap-collision calculation above. Hence every mixed and pure artificial
principal part in the double-cap expansion cancels. This remains true on
the diagonal \(s_1=s_2\), with arbitrary pole orders and coincident
content poles. Holomorphic remainder integrals and other local residue
groups do not change that conclusion.

Fixed individual high-order residues may display powers of \(\log U\)
when \eqref{res:C12} is expanded. This is not used as a global estimate:
orders may grow with height, and individual Laurent coefficients are not
controlled here. The complete contour bounds establish convergence and
power costs uniformly without a multiplicity-dependent estimate.

\subsection{\texorpdfstring{The full \(U,X,R\) powers after two moving
content
residues}{The full U,X,R powers after two moving content residues}}\label{the-full-uxr-powers-after-two-moving-content-residues}

Consider the common-\(s\) cap representation before the direct-cutoff
modification. For content poles \(v_1,v_2\), the two moving cap
locations are \[
 w_i=(z+v_i)/2-s.
\] Every term in the corresponding differentiated residue carries the
basic exponential factor \[
 U^{w_1+w_2}=U^{z+(v_1+v_2)/2-2s}.
\] Differentiation adds logarithmic factors but does not change this
basic power.

With the global outer factor \(X^sR^{-z}\) and \(U^2=2X\), the product
is exactly

\begin{equation}\label{res:C17}
X^sR^{-z}U^{z+(v_1+v_2)/2-2s}
 =2^{-s}(U/R)^zU^{(v_1+v_2)/2}.
\end{equation}

For \(R=X^\gamma\), this is

\begin{equation}\label{res:C18}
2^{-s+z/2+(v_1+v_2)/4}
 X^{(1/2-\gamma)z+(v_1+v_2)/4}.
\end{equation}

Thus shifting the \(s\)-line does not leave an uncanceled factor
\(X^{\Re s}\) that would furnish a saving. The compensating \(U^{-2s}\)
removes it.

For the actual cofactor-dependent cutoff, the per-cofactor outer factor
is \((X/n)^{s-\gamma z}\), so the exact identity becomes

\begin{equation}\label{res:C19}
(X/n)^{s-\gamma z}
 U^{z+(v_1+v_2)/2-2s}
 =
 2^{-s+z/2+(v_1+v_2)/4}
 n^{-(s-\gamma z)}
 X^{(1/2-\gamma)z+(v_1+v_2)/4}.
\end{equation}

The power of \(n\) is precisely the shifted cofactor argument already
found in \eqref{res:C3}. If the direct cutoff is Mellin-inverted with
\(h\), replace \(s\) by \(s-h\) in the cap locations; the \(X^{s-h}\)
factor cancels the corresponding \(U^{-2(s-h)}\) in the same way.

On complete block contours, \(\Re v_i\leq v_R\), and \(\Re z=\alpha\).
The resulting \(X\)-power is at most

\begin{equation}\label{res:C20}
\alpha(1/2-\gamma)+v_R/2
 =\sigma+\frac{v_R-v_L}{2}-(2\sigma-v_L)\gamma.
\end{equation}

This is exactly the direct double-residue cost in the direct
double-residue bounds. In the balanced sum, fixed powers of \(n\) and
the finite cofactor count cost \(X^\varepsilon\). No support improvement
follows by omitting the cap powers or treating a pole of an uncapped
term as a pole of the finite \(M_U\).


\subsection{The larger holomorphy domain of the numerator}

The product for \(G=\B E_\infty\) also proves holomorphy on
\[
 \Rea s>1/2,\qquad \Rea v>0,\qquad \Rea(s+2v)>\beta.
\]
Indeed, only the even shifted factors \(P(s+2jv)\) are inverted.
Their real parts are greater than \(\beta\), so they are nonzero.
Every odd shifted factor is entire, and the tail converges normally
in logarithms after entering an ordinary Euler half-plane.
The function \(E_\infty\) is holomorphic on this domain as well.
This observation does not imply that \(G\) is nonzero.

For a global cutoff, a formal content residue in the zero-cap-shift
term, at \(v_0=2s-z\), has the real-power factor
\[
 X^sR^{-2s+v_0}=(X/R^2)^sR^{v_0}.
\]
At \(\Rea s=\beta+\varepsilon\) its limiting size is
\(X^{\beta+\varepsilon}R^{-(2\beta-\Rea v_0)}\), reproducing
the old scale at the largest permitted real part \(u\).
A further leftward norm shift would have to account for the
logarithmic-derivative poles and bound their complete contribution.
This observation is not a lower bound for the finite capped sum.

The content weight is also not multiplicative in ideal convolution.
At a split prime,
\(\cont(\mathfrak p)=\cont(\overline{\mathfrak p})=1\), whereas
\(\cont(\mathfrak p\overline{\mathfrak p})=\cont((p))=p\).
Thus the ordinary reciprocal convolution identity for \(P\)
does not imply cancellation of the product of separate content
weights.

\subsection{Regularity and the unresolved cancellation}

The stated derivative condition on \(W\) permits strong
subexponential decay of its one-dimensional Mellin transform
after optimizing integration by parts. In the two-cap weight,
however, \(W\) controls the sum of the two logarithmic norm
coordinates, while their difference is controlled by the caps.
The caps are only assumed smooth. In particular, two content poles
can have large opposite ordinates, making the sum frequency small.
The stronger regularity of \(W\) alone gives no corresponding
stronger bound in that difference direction.

Finite smooth seminorms already suffice for the polynomially
bounded complete-block construction. Additional decay at large
height would not by itself supply the missing fixed power saving
for the complete contribution, including its low-height blocks.
Neither the analytic continuation barrier nor an upper-bound
exponent proves that cancellation in the actual weighted
functional is impossible.


\section{Certified arithmetic and the precise unresolved
estimate}\label{certified-arithmetic-and-the-precise-unresolved-estimate}

\subsection{Exact rational endpoint
arithmetic}\label{exact-rational-endpoint-arithmetic}

Under the stated Liu input, use the enclosure \[
 B_-=
 0.873345856106718576801184796065
 <B_*<
 0.873345856106718576801184796066
 =B_+.
\] Take \(\beta=\rho=B_*\), \(u=B_*/2\), \(v_L=1/20\), and
\(\gamma=7/25\). The enclosure gives \[
 0.063327071946640711599407601967
 <\ell<
 0.0633270719466407115994076019675,
\] so \(v_L=0.05<\ell\). Fixed sufficiently small positive \(\delta\)
and \(v_R-u\) satisfy every contour constraint, including \(\alpha<3\).

The limiting exponents are the exact linear expressions

\begin{equation}\label{res:D1}
E_I=\frac{11}{25}B_*+\frac7{500},
 \qquad
 E_S=\frac{113}{200}B_*+\frac3{2000},
\end{equation}

\begin{equation}\label{res:D2}
E_D^{\rm direct}=\frac{69}{100}B_*-\frac{11}{1000},
 \qquad
 E_D^{\rm available}=E_J=\frac{29}{50}B_*.
\end{equation}

The last identity uses \(\gamma=7/25>1/4\), so the complete-functional
recovery \eqref{res:B20} applies.

Exact rational endpoint arithmetic gives the following values, with the
strict comparisons established by the printed enclosure:

\begin{center}
\begin{tabular}{lll}
Quantity & Exact limiting exponent & Approximate value\\ \hline
\(E_I\) & \(11B_*/25+7/500\) & \(0.3982721766869562\)\\
\(E_S\) & \(113B_*/200+3/2000\) & \(0.4949404087002960\)\\
\(E_D^{\rm direct}\) & \(69B_*/100-11/1000\) & \(0.5916086407136358\)\\
\(E_D^{\rm available}\) & \(29B_*/50\) & \(0.5065405965418968\)
\end{tabular}
\end{center}

Thus both the bulk and the complete singles are square-root admissible
with a fixed positive margin before the arbitrarily small losses are
chosen. For example, the upper endpoint of \(E_S\) is
\(0.49494040870029599589266940977729\), and the lower endpoint of
\(E_D^{\rm available}\) is \(0.5065405965418967745446871817177\). The
smaller square-root margin is

\begin{equation}\label{res:D3}
\frac12-E_S>
 0.00505959129970400410733059022271.
\end{equation}

The available complete-double exponent is strictly greater than \(1/2\).

\subsection{Complete weighted
functionals}\label{complete-weighted-functionals}

Define the actual complete weighted double functional by

\begin{equation}\label{res:D4}
\mathscr D_\gamma(X)
 =\sum_{N\mathfrak c<2L_X}a(\mathfrak c)\,
 D_{V_{N\mathfrak c}}
 \left(X/N\mathfrak c,\sqrt{2X},(X/N\mathfrak c)^\gamma\right).
\end{equation}

Define \(\mathscr I_\gamma,\mathscr S_{1,\gamma},
\mathscr S_{2,\gamma},\mathscr J_\gamma\) in the same way. Uniform
profile bounds and elementary ideal counting permit summing the exact
decomposition and the bounds. For any fixed \(\varepsilon>0\), absolute
convergence of the differentiated series \(\zeta_K(1+\varepsilon)^2\)
gives \[
 \sum_{\N\mathfrak c\leq L}
 d_{2,K}(\mathfrak c)\log\N\mathfrak c
 \ll_\varepsilon L^{1+\varepsilon}.
\] Since \(L_X=X^{o(1)}\), this count and any fixed powers of \(n<2L_X\)
cost only \(X^\varepsilon\). Thus

\begin{equation}\label{res:D5}
\mathscr J_\gamma
 =\mathscr I_\gamma+\mathscr S_{1,\gamma}
  +\mathscr S_{2,\gamma}+\mathscr D_\gamma.
\end{equation}

At \(\gamma=7/25\), for every \(\varepsilon>0\),

\begin{equation}\label{res:D6}
\mathscr I_\gamma\ll_\varepsilon X^{E_I+\varepsilon},
 \qquad
 |\mathscr S_{1,\gamma}|+|\mathscr S_{2,\gamma}|
 \ll_\varepsilon X^{E_S+\varepsilon},
\end{equation}

and the strongest bound furnished by these estimates is

\begin{equation}\label{res:D7}
\boxed{\quad
 \mathscr D_{7/25}(X)
 \ll_\varepsilon X^{(29/50)B_*+\varepsilon}.
 \quad}
\end{equation}

This follows even by summing absolute values over the finite cofactor
set; it therefore also bounds the signed or complex weighted functional.
It is not a lower bound on that functional.

The old square-root content-cutoff exponent is \[
 \gamma_{\rm old}=\frac{2B_*-1}{3B_*}.
\] Using monotonicity and exact rational endpoint division gives

\begin{equation}\label{res:D8}
\begin{gathered}
0.2849927501963953104291995149000268993<\gamma_{\rm old},\\
 \gamma_{\rm old}<0.2849927501963953104291995149004639243.
\end{gathered}
\end{equation}

Thus \(7/25=0.28\) is strictly smaller than the exponent justified by
the old bound.

For comparison, the single-residue square-root threshold is \[
 \gamma_S=
 \frac{B_*+(B_*/2-1/20)/4-1/2}{2B_*-1/20},
\] and exact rational bounds put it between

\begin{equation}\label{res:D9}
0.2770179666327812371647371330282040836
 \quad\hbox{and}\quad
 0.2770179666327812371647371330285405998.
\end{equation}

The chosen exponent \(0.28\) lies above this threshold.

\subsection{The additional estimate that would suffice, and its exact
saving}\label{the-additional-estimate-that-would-suffice-and-its-exact-saving}

At \(\gamma=7/25\), the established bulk and single bounds lie strictly
below \(1/2\). Therefore \eqref{res:D5} implies the following
equivalence of square-root targets:

\begin{equation}\label{res:D10}
\mathscr J_{7/25}(X)\ll_\varepsilon X^{1/2+\varepsilon}
 \quad\Longleftrightarrow\quad
 \mathscr D_{7/25}(X)\ll_\varepsilon X^{1/2+\varepsilon}.
\end{equation}

The missing additional estimate is exactly the right-hand side for the
complete weighted double functional. It may exploit cancellation across
blocks or cofactors. It need not, and should not without proof, be
replaced by a demand for absolute convergence or a square-root bound on
every individual zero residue.

Relative to \eqref{res:D7}, the necessary exponent improvement is \[
 \vartheta_*=\frac{29}{50}B_*-\frac12,
\] with exact strict enclosure

\begin{equation}\label{res:D11}
0.0065405965418967745446871817177
 <\vartheta_*<
 0.00654059654189677454468718171828.
\end{equation}

A bound \[
 \mathscr D_{7/25}(X)
 \ll_\varepsilon X^{(29/50)B_*-\vartheta+\varepsilon}
\] with \(\vartheta\geq\vartheta_*\) supplies the square-root target;
\(\vartheta>\vartheta_*\) supplies a fixed power margin. The reduction
in the content-cutoff exponent would be

\begin{equation}\label{res:D12}
\gamma_{\rm old}-7/25
 =\frac{\vartheta_*}{3B_*/2}
 \approx0.004992750196395310429199514900.
\end{equation}

More generally, taking \(\gamma=\gamma_{\rm old}-d\geq1/4\), where
\(d>0\) and the already established single bound remains square-root
admissible, requires an improvement of \((2\beta-u)d\) over the
inherited old exponent for the complete double functional. The existence
of such an improvement is not proved by this calculation.

\section*{Method, record, and dedication}

The local extraction, convergence, fixed-parameter nonvanishing,
and multiplicity-aware cancellation count are proved in full. The complete grouped-residue decomposition, its convergence and bounds, the cofactor identities, and arbitrary-order cap cancellations are also proved in full. The numerical residue example uses Liu's stated zero-free theorem as an external input, without reconstructing its analytic proof or adding zeta RH.
The published zero-count theorems and Hecke continuation and
strip-growth results are the analytic inputs specified above.
The reciprocal bound is proved here by a positive harmonic-function
argument. The proofs of the two published counting theorems have
not been independently reconstructed.

The mathematics and text were produced by GPT-6 Astra in ChatGPT
Pro mode. Separate instances of the same model re-proved the main
statements from their statements alone and performed targeted
source and proof checks. Their reasoning effort was inherited and
not exposed. No human has checked this work. The exact check scopes,
programs, certificates, and source locators are recorded in the
\href{\workbenchurl}{workbench repository} and the
\href{\recordurl}{permanent research record}.
The original text is dedicated to the public domain under
\href{https://creativecommons.org/publicdomain/zero/1.0/}{CC0 1.0}.
Cited publications retain their respective rights.

\begin{thebibliography}{9}
\bibitem[Bennett et al.(2021)]{BMOR}
Bennett, M. A., Martin, G., O'Bryant, K., \& Rechnitzer, A. (2021).
Counting zeros of Dirichlet $L$-functions.
\emph{Mathematics of Computation, 90}(329), 1455--1482.
\url{https://doi.org/10.1090/mcom/3599}.

\bibitem[Hasanalizade et al.(2022)]{HSW}
Hasanalizade, E., Shen, Q., \& Wong, P.-J. (2022).
Counting zeros of the Riemann zeta function.
\emph{Journal of Number Theory, 235}, 219--241.
\url{https://doi.org/10.1016/j.jnt.2021.06.032}.

\bibitem[Liu(2026)]{AMLiu}
Liu, B. (2026).
\emph{Joint marked moments and optimized zero-free half-planes for
unitary Hecke L-functions} (Version~2). Hexagon, 2611.00421v2.
\url{https://hexagonmath.org/2611.00421v2}.
Theorem~1.2, p.~3; equations~(42), (44), pp.~27--29;
Theorem~8.29, pp.~78--80. The zero-free theorem is imported as a
stated external input; its analytic proof is not reconstructed here.

\bibitem[Zaman(2016)]{Zaman}
Zaman, A. (2016).
Explicit estimates for the zeros of Hecke $L$-functions.
\emph{Journal of Number Theory, 162}, 312--375.
\url{https://doi.org/10.1016/j.jnt.2015.10.003}.
Primary preprint: \url{https://arxiv.org/pdf/1502.05679},
version 4, \S2.1, equations (2.1)--(2.3), Lemmas 2.1--2.3,
preprint pp.~6--7. Cited only for the completion, functional equation,
and fixed-strip growth background, not for a unit-height zero count
or the reciprocal bound proved here.
\end{thebibliography}
\end{document}
